\chapter*{Lista de Símbolos}
\addcontentsline{toc}{chapter}{Lista de Símbolos}

% PLACEHOLDER: preencher apenas se o corpo do texto usar notacao matematica
% formal (ex.: metrica de erro de sincronizacao musica-evento no Cap03, ou
% modelo linear de efeitos mistos na analise estatistica). Se o TCC nao usar
% simbolos formais suficientes, considerar remover esta lista e o comando
% \simbolos em main.tex (combinar com o orientador).

\begin{description}
    \item[$\overline{e}$] Erro médio de sincronização entre eventos do jogo e batidas da música
    \item[$e_i$] Erro de sincronização do $i$-ésimo evento em relação à batida mais próxima
    \item[$\Delta$] Variação de uma medida entre dois momentos (ex.: pós-estressor e pós-intervenção)
    \item[$\ln(\mathrm{RMSSD})$] Logaritmo natural da RMSSD, empregado porque a distribuição da VFC costuma ser assimétrica
    % PLACEHOLDER: adicionar os simbolos do modelo estatistico (Cap03/Cap04),
    % ex.: coeficientes de efeitos fixos e aleatorios do modelo linear misto.
\end{description}
